% =========================================================
% Proof: Completion Theorem
% Source: notes/completeness/notes-completeness.tex
% Difficulty: Hard — defer until comfortable with equivalence classes
% =========================================================

\subsection*{Completion Theorem}
\label{prf:completion}

% ⚠️  DIFFICULTY WARNING
% This proof is long and structurally complex.
% It requires:
%   - Equivalence classes of Cauchy sequences (co-Cauchy relation)
%   - Defining a metric on equivalence classes (well-definedness check)
%   - Embedding the original space isometrically
%   - Proving the completion is complete (diagonal argument on double sequences)
%   - Proving the embedding is dense
% Recommended: return to this proof after the compactness chapter,
% when you are fluent with the co-Cauchy construction and diagonal arguments.

\begin{remark*}[Return]
$\leftarrow$ \hyperref[thm:completion]{Back to Theorem in Notes}
\end{remark*}

\begin{proof}
\Claim Every metric space $(X,d)$ has a completion: a complete metric space
$(\hat{X}, \hat{d})$ together with an isometric embedding $\iota : X \to \hat{X}$
such that $\iota(X)$ is dense in $\hat{X}$.
Moreover, the completion is unique up to isometry.

\Given A metric space $(X,d)$.

\Goal To construct $(\hat{X}, \hat{d})$ and $\iota$, and verify:
\begin{enumerate}[(i)]
  \item $\hat{d}$ is a well-defined metric on $\hat{X}$.
  \item $\iota$ is an isometric embedding.
  \item $\iota(X)$ is dense in $\hat{X}$.
  \item $(\hat{X}, \hat{d})$ is complete.
\end{enumerate}

\Strategy

\noindent\textbf{Construction.}
Let $\hat{X}$ be the set of equivalence classes of Cauchy sequences in $X$
under the co-Cauchy relation $(x_n) \sim (y_n) \iff d(x_n, y_n) \to 0$
(Proposition~\ref{prop:co-cauchy-equivalence-relation}).

Define $\hat{d}([(x_n)], [(y_n)]) := \lim_{n \to \infty} d(x_n, y_n)$.

\noindent\textbf{Step 1: Well-definedness of $\hat{d}$.}
Show $\lim d(x_n, y_n)$ exists (the sequence is real and Cauchy by reverse
triangle inequality) and is independent of the representatives chosen.

\noindent\textbf{Step 2: $\hat{d}$ is a metric.}
Verify M1--M4 for $\hat{d}$, using the corresponding properties of $d$
and limit arithmetic.

\noindent\textbf{Step 3: Embedding.}
Define $\iota(x) = [(x, x, x, \ldots)]$ (constant sequence).
Show $\hat{d}(\iota(x), \iota(y)) = d(x,y)$.

\noindent\textbf{Step 4: Density.}
Show every $[(x_n)] \in \hat{X}$ is the limit of $(\iota(x_n))_{n}$.
Use $\hat{d}(\iota(x_n), [(x_k)]_k) = \lim_k d(x_n, x_k) \to 0$ as $n \to \infty$.

\noindent\textbf{Step 5: Completeness.}
Take a Cauchy sequence $([(x_n^{(j)})]_j)$ in $\hat{X}$;
use density to pick representatives in $\iota(X)$;
extract a diagonal Cauchy sequence in $X$;
show its equivalence class is the limit.

% -------------------------------------------------------
% YOUR PROOF HERE (work through each step above)
% -------------------------------------------------------

\end{proof}

\begin{remark*}[Proof shape]
Construction proof in five stages: define $\hat{X}$, verify $\hat{d}$ is
well-defined and is a metric, embed $X$ isometrically, show the image is dense,
prove completeness via a diagonal argument.
Each step is independently checkable.
\end{remark*}

\begin{remark*}[Dependencies]
\begin{itemize}
  \item Proposition~\ref{prop:co-cauchy-equivalence-relation}: co-Cauchy is an equivalence relation.
  \item Proposition~\ref{prop:cauchy-bounded}: Cauchy sequences are bounded.
  \item Proposition~\ref{prop:reverse-triangle}: reverse triangle inequality
        (to show $d(x_n, y_n)$ is Cauchy in $\mathbb{R}$).
  \item Proposition~\ref{prop:convergent-implies-cauchy}: convergent $\Rightarrow$ Cauchy
        (for Step 4).
  \item Completeness of $\mathbb{R}$ (limits of real Cauchy sequences exist).
  \item Definition~\ref{def:isometry}: isometric embedding.
\end{itemize}
\end{remark*}
